\documentclass[12pt]{article}
\usepackage{amsmath}
\usepackage{amssymb}
\usepackage{booktabs}

%--------------------------------------
\newcommand\qgb{Q_{\rm GB}}
\newcommand\qgbj{Q_{{\rm GB},j}}
\newcommand\vpr{V^\prime}
\newcommand\qhat{\widehat{Q}}
\newcommand{\sj}{{\sigma,j}}
%--------------------------------------

\begin{document}
\begin{center}
{\bf EGYRO: formulation}\\
{\sl J. Candy}
\end{center}

\section{Transport problem}

Steady state, for each energy channel $\sigma \in \{e,i\}$:
%
\begin{equation}
\frac{1}{\vpr(r)} \frac{\partial}{\partial r}
 \left[ \vpr(r) \, Q_\sigma(r) \right] = S_\sigma(r) \; .
\label{eq.transport}
\end{equation}
%
Integrating gives the target flux
%
\begin{equation}
Q^T_\sigma(r) = \frac{P_\sigma(r)}{\vpr(r)} \; ,
\qquad
P_\sigma(r) = \int_0^r dx \, \vpr(x) \, S_\sigma(x) \; ,
\label{eq.target}
\end{equation}
%
so only the integrated powers enter.  The sources depend on the
profiles, $S_\sigma = S_\sigma(T)$, through fusion heating,
electron-ion exchange and radiation.  The transport flux is
%
\begin{equation}
Q_\sigma = Q_\sigma^{\rm turb} + Q_\sigma^{\rm neo} \; ,
\end{equation}
%
with $Q^{\rm turb}_\sigma$ from TGLF and, at present,
$Q^{\rm neo}_\sigma = 0$.

\section{Normalization}

\begin{equation}
\qgb = n_e T_e c_s \left( \frac{\rho_s}{a} \right)^{2} \; ,
\qquad
c_s = \sqrt{\frac{T_e}{m_D}} \; ,
\qquad
\rho_s = \frac{m_D c_s}{e B_{\rm unit}} \; ,
\end{equation}
%
with $a$ the minor radius of the last closed flux surface.  Explicitly,
%
\begin{equation}
\qgb = \frac{m_D^{1/2}}{e^2 a^2 B_{\rm unit}^2} \; n_e \, T_e^{5/2} \; .
\label{eq.qgb}
\end{equation}
%
Writing $\qhat \doteq Q/\qgb$, the normalized target is
%
\begin{equation}
\qhat^T_\sigma = \frac{P_\sigma}{\vpr \, \qgb} \; ,
\label{eq.qhattarget}
\end{equation}
%
which depends on the solution through $\qgb$ even when $S_\sigma$ does
not.

\section{Scale-length formulation}

With $x \doteq r/a$, the unknown is
%
\begin{equation}
z_\sigma(x) \doteq - \frac{\partial \ln T_\sigma}{\partial x} \; .
\label{eq.zdef}
\end{equation}
%
Integrating (\ref{eq.zdef}) determines the profile {\it shape}
%
\begin{equation}
\Phi_\sigma(x;z) = \exp \left( \int_x^{x_{\rm ref}} dx^\prime \,
 z_\sigma(x^\prime) \right) \; ,
\qquad
T_\sigma \propto \Phi_\sigma \; ,
\label{eq.shape}
\end{equation}
%
fixed by $z_\sigma$ alone.  The reference radius $x_{\rm ref}$ is
arbitrary, so $\Phi_\sigma$ is defined only up to a multiplicative
constant.  Note $\Phi_\sigma > 0$ for every $z_\sigma$: positivity of
$T_\sigma$ is automatic.

\subsection{Axis condition}

$T_\sigma$ is even in $x$, so $z_\sigma$ is odd:
%
\begin{equation}
z_\sigma(0) = 0 \; .
\label{eq.axis}
\end{equation}
%
This is imposed, not solved for.

\subsection{Grid and reintegration}

\begin{equation}
x_0 = 0 \; , \qquad x_1 < x_2 < \dots < x_N \; ,
\qquad x_1 \ge 0.2 \; ,
\label{eq.grid}
\end{equation}
%
non-uniform, with resolution concentrated toward the edge.  The
integral in (\ref{eq.shape}) is a linear functional of $z_\sigma$, so
for any quadrature with weights $W$,
%
\begin{equation}
\boxed{\;
\ln \Phi_{\sj} = \sum_{m=0}^{N} W_{jm} \, z_{\sigma,m} \; ,
\qquad z_{\sigma,0} = 0 \; . }
\label{eq.W}
\end{equation}
%
$W$ is known, built once from the grid, and requires no flux
evaluation.  Positivity holds for any $W$.

\section{Boundary condition}
\label{sec.bc}

Any {\it linear} condition on the discrete profile $T_\sigma = \{
T_{\sj} \}$ may be written, without loss of generality, as
%
\begin{equation}
B_\sigma \cdot T_\sigma = 1 \; ,
\label{eq.linbc}
\end{equation}
%
with $B_\sigma$ a known vector.  Since $T_\sigma \propto \Phi_\sigma$,
this fixes the constant of proportionality, and the profile follows
from $z$ alone:
%
\begin{equation}
\boxed{\;
T_\sigma = \frac{\Phi_\sigma(z)}{B_\sigma \cdot \Phi_\sigma(z)} \; }
\label{eq.profile}
\end{equation}
%
The right side is homogeneous of degree zero in $\Phi_\sigma$, so the
reference radius $x_{\rm ref}$ cancels identically.

\section{Discrete problem}

With $n_c$ channels the unknown vector is
%
\begin{equation}
z = \{ z_\sj \} \in \mathbb{R}^{N n_c} \; ,
\qquad
z_\sj \doteq z_\sigma(x_j) \; ,
\end{equation}
%
and the $N n_c$ equations to be solved are
%
\begin{equation}
\boxed{\;
\qhat_\sj(z) = \qhat^T_\sj(z) \; }
\label{eq.match}
\end{equation}
%
both sides evaluated on the profile (\ref{eq.profile}) built from the
same $z$.  The measure of departure from (\ref{eq.match}) is a property
of the solution method, and is not specified here.

\section{Structure}

The two sides of (\ref{eq.match}) have opposite structure.

\noindent {\bf The flux is local.}  $\qhat_\sj$ depends on $z$ through
the $n_c$ components at $x_j$ alone, and on the profile through two
scalars:
%
\begin{equation}
\qhat_\sj = \qhat_\sigma \bigl( z_{\cdot,j} ;\; T_{e,j},\, T_{i,j}
 \bigr) \; ,
\label{eq.localflux}
\end{equation}
%
with density, geometry and field fixed.  The profile enters through
$\tau = T_i/T_e$, $\beta_e \propto n_e T_e/B_{\rm unit}^2$, $\nu_e
\propto n_e T_e^{-2}$, the Debye length, and $\qgb \propto n_e
T_e^{5/2}$.  This dependence is weak --- hence the gyroBohm
normalization --- but nonzero.

\noindent {\bf The target is nonlocal.}  By (\ref{eq.target}),
%
\begin{equation}
\qhat^T_\sj = \frac{1}{\vpr_j \, \qgbj}
 \int_0^{x_j} dx \, \vpr(x) \, S_\sigma(x) \; ,
\label{eq.nonlocaltarget}
\end{equation}
%
an integral over the volume enclosed by $x_j$.  Since $S_\sigma$
depends on the local profile, $\qhat^T_\sj$ is a functional of
$T_\sigma$ at every radius {\it interior} to $x_j$, and through
$\qgbj$ of $T_{e,j}$.

\noindent {\bf The profile is nonlocal in $z$.}  By (\ref{eq.profile}),
$T_\sj$ depends on $z_\sigma$ at every radius: through $W$, and through
the normalization $B_\sigma \cdot \Phi_\sigma$.

\noindent Hence (\ref{eq.match}) is coupled across the grid.  Were the
profile dependence of both sides absent, it would separate into $N$
independent problems of dimension $n_c$.

\noindent {\bf Stiffness.}  $\qhat_\sigma$ is small below a critical
gradient and rises steeply above it.  Recovering $z$ from a given flux
is therefore well conditioned above threshold, and ill-posed below it,
where a root may be non-unique or absent.

\section{Solver}

EGYRO evaluates $\qhat_\sj(z)$ and $\qhat^T_\sj(z)$ and delegates the
solve to a global, derivative-free method.  Relevant structure:
%
\begin{itemize}
\item Every $z$ is admissible, so $0 \le z_\sj \le z_{\rm max}$ is a box
domain with no constraints to enforce.
\item One evaluation of (\ref{eq.match}) is $N$ independent flux
calculations, dispatched concurrently: cost is measured in sequential
batches, not flux calls.
\item By (\ref{eq.localflux}), one evaluation samples the same local
flux function at $N$ points --- $N$ training points per evaluation for
a surrogate.
\item Holding $T_\sj$ fixed separates (\ref{eq.match}) into $N$
problems of dimension $n_c$; reintegrating then moves $T$ and the
process repeats.  This is exact at the fixed point, and its cost is the
outer iteration count, set by the temperature dependence above.
\item Scaling $P_\sigma \to \epsilon P_\sigma$ and continuing $\epsilon:
0 \to 1$ traces a path from the zero-flux state to the solution.
\end{itemize}

\section{Implementation}

\begin{center}
\begin{tabular}{ll}
\toprule
Module & Role \\
\midrule
\verb|api.py|                & flux and target (\ref{eq.match}); $z \to$ profile map \\
\verb|transport/profiles.py| & profile load; grid; reintegration (\ref{eq.W}) \\
\verb|transport/tglf_map.py| & profile state $\to$ TGLF inputs \\
\verb|transport/fluxes.py|   & concurrent flux evaluation; $\qhat_\sj$ \\
\verb|transport/sources.py|  & target $\qhat^T_\sj$, Eq.~(\ref{eq.qhattarget}) \\
\verb|transport/output.py|   & diagnostic output \\
\bottomrule
\end{tabular}
\end{center}

\section{Current restrictions}

\begin{enumerate}
\item $S_e$, $S_i$ are evaluated once from the experimental profiles
rather than at each iteration --- a convenience for method development,
not a property of the formulation.  The code therefore misses the
nonlocality of (\ref{eq.nonlocaltarget}), retaining only the dependence
of $\qhat^T_\sj$ on $T_{e,j}$ through $\qgbj$.
\item $Q_\sigma^{\rm neo} = 0$.
\item Energy channels only: density gradients are held fixed; no
particle or momentum channel.
\item Each ion species carries an independent gradient against a single
ion target, split evenly.  The physical formulation is one ion
temperature with the summed ion flux matched to $Q^T_i$, $n_c = 2$.
\item Only the Dirichlet condition at $x_N$ is implemented;
Sec.~\ref{sec.bc} is not.
\item $W$ is the trapezoidal rule on a uniform grid: second order, and
in error by $26\%$ in $T_e$ at $N=4$ ($5\%$ at $N=8$), amplified by
$5/2$ in $\qhat^T$.
\item Eq.~(\ref{eq.axis}) is not enforced; the interpolated
experimental value is carried at $x=0$.
\item $q$, $s$, $B_{\rm unit}$ and $\partial \ln p_{\rm tot}/\partial r$
are not updated as the profiles evolve.
\label{item.frozen}
\end{enumerate}

\section{A learned inverse}

Existing surrogate methods (PORTALS) fit a model of the {\it forward}
map (\ref{eq.localflux}), $z \mapsto \qhat$, and then solve
%
\begin{equation}
\qhat^{\rm model}_\sj (z) = \qhat^T_\sj
\end{equation}
%
for $z$ by numerical root-finding, at every iteration.  We propose to
fit the {\it inverse} map instead.

\subsection{The inverse map}

At fixed $T_{\cdot,j}$, (\ref{eq.localflux}) is $n_c$ equations in the
$n_c$ unknowns $z_{\cdot,j}$, and may be inverted:
%
\begin{equation}
\boxed{\;
z_{\cdot,j} = \zeta \bigl( \qhat_{\cdot,j},\, T_{\cdot,j} ;\, g_j
 \bigr) \; }
\label{eq.inverse}
\end{equation}
%
$\zeta$ returns the gradient required to carry a {\it given} flux.
Since the target $\qhat^T$ is known and cheap to compute, the gradient
that matches it is obtained by {\it evaluating} $\zeta$ --- no
root-finding is performed at any stage.

\subsection{Why the inverse is the easier function}

By Sec.~6 the flux is stiff: $\qhat_\sigma$ spans orders of magnitude
over a narrow range of $z$.  Two consequences, in opposite directions:
%
\begin{itemize}
\item $\qhat(z)$ is steep, hence expensive to fit accurately, and a
small error in the fitted flux is {\it amplified} when it is inverted
to recover $z$.
\item $\zeta(\qhat)$ is flat, hence cheap to fit accurately, and an
error in the fitted gradient is {\it not} amplified --- it is the
answer.
\end{itemize}
%
The existing approach fits the hard direction and then does work to
invert it.  Eq.~(\ref{eq.inverse}) fits the easy direction and reads
the answer off.

The training data are identical.  Each flux evaluation at $x_j$ yields
a pair $(z_{\cdot,j}, T_{\cdot,j}) \to \qhat_{\cdot,j}$.  For a forward
model this is input $(z,T)$, output $\qhat$; for the inverse it is
input $(\qhat,T)$, output $z$.  The same samples, relabelled, at no
additional cost.

\subsection{The iteration}

With $\zeta$ in hand, flux matching is a fixed point in the profile
alone:
%
\begin{equation}
z^{(k+1)}_{\cdot,j} = \zeta \bigl( \qhat^T_{\cdot,j}(z^{(k)}),\,
 T_{\cdot,j}(z^{(k)}) ;\, g_j \bigr) \; ,
\label{eq.picard}
\end{equation}
%
with $T(z)$ from (\ref{eq.W}) and (\ref{eq.profile}).  Each pass
evaluates the target, evaluates $\zeta$ at every radius, and
reintegrates.  Were $T$ exact, one pass would suffice; the residual
inconsistency is only the {\it weak} temperature dependence of Sec.~6,
which sets the rate.  There is no optimization in $N n_c$ dimensions,
no acquisition function, and no Jacobian.

\subsection{Amortization}

At fixed $g_j$, $\zeta$ is a fixed function of $2 n_c$ arguments, the
same for every discharge.  Since TGLF is cheap, it may be trained
offline over the dimensionless input space once and for all, rather
than refitted per radius and per discharge from the samples of a single
run.

\subsection{Open}

$\zeta$ is single-valued only above the critical gradient; below it
$\qhat_\sigma \simeq 0$, the flux carries no information about the
gradient, and the inverse is ill-posed (Sec.~6).  The contraction rate
of (\ref{eq.picard}) is unmeasured.

\section{Edge boundary condition}

EGYRO solves to the separatrix.  ELM-free and without a pedestal, there
is no pedestal top at which to pin $T_\sigma$, and the edge temperature
is set by parallel transport in the scrape-off layer,
%
\begin{equation}
T_{\sigma,N} = f_\sigma \bigl( P_{\rm sep} \bigr) \; ,
\qquad
T_{\sigma,N} \sim P_{\rm sep}^{2/7} \; ,
\label{eq.sol}
\end{equation}
%
the exponent following from $q_\parallel = -\kappa_0 T^{5/2} \partial_s
T$.  But the power crossing the separatrix,
%
\begin{equation}
P_{\rm sep} = \int_0^1 dx \, \vpr(x) \, S_\sigma(x) \; ,
\end{equation}
%
is a functional of the profile, since $S_\sigma = S_\sigma(T)$.  The
edge amplitude and the integrated source therefore close on each other:
$B_\sigma$ in (\ref{eq.linbc}) is itself a function of $z$, and the
closed form (\ref{eq.profile}) is unavailable.

Take $x_{\rm ref} = x_N$, so that $A_\sigma = T_{\sigma,N}$ and
$T_\sigma = A_\sigma \Phi_\sigma(z)$.  Promote $A_\sigma$ to an unknown
and append (\ref{eq.sol}):
%
\begin{align}
\qhat_\sj(z,A) = &~ \qhat^T_\sj(z,A) \; ,
 && N n_c \ \text{equations} \; , \\
A_\sigma - f_\sigma \bigl( P_{\rm sep}[A_\sigma \Phi_\sigma(z)] \bigr)
 = &~ 0 \; , && n_c \ \text{equations} \; ,
\label{eq.augmented}
\end{align}
%
with $(z,A) \in \mathbb{R}^{N n_c + n_c}$.  The appended rows require no
flux evaluation.  The $2/7$ exponent makes $T_{\sigma,N}$ respond weakly
to $P_{\rm sep}$, so the added nonlinearity is soft.

\end{document}
